% kernel/engineering_decisions-DE.tex

\chapter{Vom bewerteten Vorschlag zum autorisierten Zustand}
\label{chap:engineering-decisions}

Die Berlinish-Komposition hat die erforderlichen Prüfungen von Geometrie,
Realisierung, Geschwindigkeit, Überhöhung und Evidenz bestanden. Sie ist nun
ein für den angegebenen Zweck anwendbarer Kandidat. Sie ist noch immer nicht
der angenommene Zustand des Alignments.

\section{Vorschlag, Kandidat und Bewertung}

\begin{definition}[Vorschlag]
\label{def:proposal}
Ein Vorschlag ist eine mögliche Ingenieuränderung, die vor Annahme oder
Ablehnung zur Bewertung eingereicht wird.
\end{definition}

\begin{definition}[Lösungskandidat]
\label{def:candidate-solution}
Ein Lösungskandidat ist ein vorgeschlagener Engineering State oder eine
vorgeschlagene Konfiguration, die anhand anwendbaren Wissens,
Nebenbedingungen und Präferenzen bewertet werden kann.
\end{definition}

axtranNew sucht oder löst im Raum ausdrückbarer Übergänge. Das
Engineering-Object-Modell trägt den entstandenen Vorschlag und seine Beziehung
zum Alignment. Bewertung bestimmt zweckrelative Eignung. Dies sind getrennte
Verantwortungen: Erfolgreiche Erzeugung bedeutet nicht Zulässigkeit, und
Zulässigkeit wählt nicht unter allen zulässigen Alternativen aus.

\section{Autorität verändert den Ingenieurstatus}

Der aktive Kernel fixiert die Grenze, dass Bewertung keine
Entscheidungsautorität ausübt; er etabliert noch keine vollständige kanonische
Identität jeder Engineering Decision. Für diese Erklärungsfolge genügt die
folgende enge operationale Formulierung:

\begin{definition}[Engineering Decision]
\label{def:engineering-decision}
Eine Engineering Decision ist eine autorisierte Annahme, Ablehnung oder
Auswahl, die einen Ingenieurvorschlag oder eine Alternative innerhalb
ausdrücklich angegebener Qualifikationen entscheidet.
\end{definition}

Die befugte Autorität kann den Kandidaten als Nachfolgezustand annehmen,
ablehnen, weitere Evidenz verlangen oder einen anderen Kandidaten auswählen.
Die Entscheidung benennt Geltungsbereich, Gegenstand, Evidenz, anwendbares
Wissen, verantwortliche Autorität und Wirkung. Sie löscht weder Vorgänger noch
verworfene Kandidaten, Vermessung oder Berechnungsprotokoll.

Bei Annahme bleibt der neue Zustand mit demselben Alignment verbunden, wenn das
autorisierte Identitätsurteil Objektkontinuität wahrt. Wird die Änderung als
Ersatzobjekt geführt, muss stattdessen diese Beziehung angegeben werden. Der
Kernel liefert keinen allgemeinen geometrischen Schwellenwert zwischen diesen
Ergebnissen; der Solver kann daher keinen erzeugen.

\section{Rückkehr zum Übergang}

Das laufende Beispiel hat die vollständige Bedeutungskette durchlaufen:

\begin{enumerate}
\item Berlinish drückte eine geordnete Übergangskonstruktion aus.
\item axtranNew oder ein gleichwertiger Dienst berechnete einen
  parametrisierten Vorschlag.
\item Der Vorschlag wurde auf ein identifizierbares Alignment und einen
  möglichen Zustand bezogen, während seine konstruktive Zusammensetzung
  erhalten blieb.
\item Metric Realization und Realization Context machten ihn räumlich messbar.
\item Repräsentationen kommunizierten ihn, und Provenienz trennte Berechnung,
  Import, Ableitung und Beobachtung.
\item Anwendbares Wissen und zweckrelative Bewertung prüften seine Eignung.
\item Eine befugte Autorität nahm, sofern sie sich dafür entschied, den
  Nachfolgezustand an.
\end{enumerate}

Zu keinem Zeitpunkt erzeugte Berechnung Autorität, erwarb Repräsentation
Identität, wurde Beobachtung zur physischen Realität oder bildete Bewertung
eine Entscheidung. Diese Ingenieurbedeutung wird nun in Geometrie und
Krümmung weitergetragen: Die Mathematik arbeitet dort mit identifizierbaren,
qualifizierten Zuständen statt mit anonymen Kurven.
